\chapter{Entanglement-Breaking Quantum Channels}
\label{chap:4_entanglement}




\subsubsection{Basics}

\begin{definition}[Entanglement breaking channel \cite{Horodecki2003}]\label{def:eb}
    Let $T \in \m C(\C^{d_A},\C^{d_B})$ be a quantum channel. It is entanglement-breaking if for all dimensions $d\in \N$ and states $\rho \in \m D(\C^d \otimes \C^{d_A})$ it holds true that $(\id_d\otimes T)(\rho)$ is separable. $T$ is said to have Holevo or measure-and-prepare form if there are density matrices $\sigma_k \in \m D(\C^{d_B})$ and a POVM $\{M_k\}_k$ on $\C^{d_A}$ so that $T(\cdot) = \sum_k \trr{M_k \cdot }\sigma_k$. The set of entanglement-breaking quantum channels $T\in \m C(\C^{d_A},\C^{d_B})$ will be denoted by $\text{EB}(\C^{d_A},\C^{d_B})$, $\text{EB}_d \coloneqq \text{EB}(\C^d)\coloneqq \text{EB}(\C^{d},\C^{d})$.
\end{definition}\noindent
Again, the order of the tensor product does not matter in this definition. By using the spectral decomposition, the density matrices and POVM effects can always be chosen to be rank-one (although the effect operators of the POVM may be multiples of each other.)
A closely related concept is that of an entanglement-annihilating channel, which destroys any entanglement within the system it acts on:
\begin{definition}[Entanglement-annihilating channel \cite{moravcikovaEntanglementannihilatingEntanglementbreakingChannels2010}]
    A quantum channel $\m E \in \m C_d$ is entanglement-annihilating if $\m E(\rho) $ is separable for all states $\rho \in \m D(\C^d)$.
\end{definition}
Clearly, a channel $\m E$ is entanglement-breaking if and only if $\id \otimes \m E$ is entanglement-annihilating.

An obvious example of a channel that is not entanglement-breaking is the identity channel $T=\id: \rho \mapsto \rho$ since then $(\id \otimes T)(\rho)$ is entangled whenever $\rho$ is.

\begin{lemma}
    A channel $T$ in Holevo form $T(\cdot) = \sum_k \trr{M_k \cdot }\rho_k$ has Kraus operators $K_{kmn}=\sqrt{\sigma_k}\ket m \bra n\sqrt{M_k}$ and Choi matrix $C_T = \sum_k M_k^\text T\otimes \sigma_k$. 
\end{lemma}
\begin{proof}
    \begin{align*}
\sum_{kmn} A_{kmn} \rho A_{kmn}^* &= \sum_{kmn} \sqrt{\sigma_k} |m\rangle\langle n| \sqrt{M_k} \rho \sqrt{M_k} |n\rangle\langle m| \sqrt{\sigma_k} \\
&= \sum_{km} \sqrt{\sigma_k} |m\rangle \text{Tr}[M_k \rho] \langle m| \sqrt{\sigma_k} = \sum_k \text{Tr}[M_k \rho] \sigma_k
\end{align*}
\end{proof}

The main result from \cite{Horodecki2003} is the following characterization of entanglement-breaking channels:
\begin{theorem}[\cite{Horodecki2003}, theorem 4]
    Let $\m E \in \m C(\C^{d_A},\C^{d_B})$ be a quantum channel, then the following statements are equivalent:
    \begin{enumerate}
        \item $\m E$ has Holevo form.
        \item $\m E$ is entanglement-breaking.
        \item The Choi matrix $C_{\m E}$ is separable.
        \item All Kraus operators of $\m E$ are rank-one.
        \item For all positive maps $\gamma$,  $\gamma \circ \mathcal{E}$ is completely positive.
        \item For all positive maps $\gamma$,  $\m E\circ\gamma$ is completely positive.
    \end{enumerate}
\end{theorem}
\begin{proof}
i) $\Rightarrow$ ii):
\begin{align*}
    (\text{id} \otimes \mathcal{E})(\rho) &= \sum_k (\text{id} \otimes \sigma_k) ( \text{id} \otimes \text{Tr}(M_k \cdot) ) (\rho) = \sum_k \text{tr}_2 \left[ (\text{id} \otimes M_k) \rho \right] \otimes \sigma_k \\
    &= \sum_k \gamma_k A_k \otimes \sigma_k,
\end{align*}
where $\gamma_k := \text{tr}[M_k \rho] \geq 0$, $\sum_k \gamma_k = 1$, and $A_k := \text{tr}_2 \left[ (\text{id} \otimes M_k) \rho \right] / \gamma_k  = \text{tr}_2[\sqrt{M_k} \rho \sqrt{M_k}]/\gamma_k\geq 0$,
implying that $(\text{id} \otimes \mathcal{E})\rho$ is separable.\\
ii) $\Rightarrow$ iii): trivial
iii) $\Rightarrow$ i):
\begin{align*}
    (\text{id} \otimes \mathcal{E})(\ket{\Omega}\bra{\Omega}) &= \sum_{ij} (\text{id} \otimes \mathcal{E})(\ket{i}\bra{j} \otimes \ket{i}\bra{j}) \\
    &= \sum_{ij} \ket{i}\bra{j} \otimes \mathcal{E}(\ket{i}\bra{j}) \\
    &\overset{3)}{=} \sum_n \tilde{p}_n \sigma_n^A \otimes \sigma_n^B = \sum_n p_n \ket{v_n}\bra{v_n} \otimes \ket{w_n}\bra{w_n}
\end{align*}
for suitably chosen normalized $\ket{v_n}, \ket{w_n}$ and $\sum p_n = 1$. Consequently, one has
\begin{align*}
    T(\rho) &= \text{tr}_1[C (\rho^T \otimes \1)] = \sum_n p_n \text{tr}_1 \left[ \ket{v_n}\bra{v_n} \rho^T \otimes \ket{w_n}\bra{w_n} \right] \\
    &= \sum_n p_n \bra{v_n} \rho^T \ket{v_n} \ket{w_n}\bra{w_n} = \sum_n p_n \bra{\bar{v}_n} \rho \ket{\bar{v}_n} \ket{w_n}\bra{w_n} \\
    &= \sum_n \text{tr} \left[ \rho \, p_n \ket{\bar{v}_n}\bra{\bar{v}_n} \right] \ket{w_n}\bra{w_n}.
\end{align*}
Now, i) - iii) are equivalent. iii) $\Rightarrow$ iv): Since
\begin{align*}
    \mathcal{E}(\rho) &= \sum_n \ket{w_n}\bra{w_n} \text{Tr}[p_n \rho \ket{v_n}\bra{v_n}] = \sum_n p_n \ket{w_n}\bra{w_n} \bra{v_n} \rho \ket{v_n} \\
    &= \sum_n p_n \ket{w_n}\bra{v_n}\rho \ket{v_n}\bra{w_n},
\end{align*}
the Kraus operators $K_n \coloneqq \sqrt{p_n}\ket{w_n}\bra{v_n}$ are rank-one.\\
iv) $\Rightarrow$ i):
\begin{align*}
    \mathcal{E}(\rho) &= \sum_k A_k \rho A_k^* = \sum_k \gamma_k \ket{v_k}\bra{w_k} \rho \ket{w_k}\bra{v_k} = \sum_k \text{Tr}[\gamma_k \rho \ket{w_k}\bra{w_k}] \ket{v_k}\bra{v_k}
\end{align*}
ii) $\Leftrightarrow$ v): Theorem \ref{thm:positive-maps} states that $\rho \text{ separable } \Leftrightarrow \forall \text{ positive } \gamma : (\text{id} \otimes \gamma)(\rho) \ge 0$ holds true for all states $\rho$. Now, one only needs to use that a channel is entanglement-breaking if and only if its Choi matrix is separable.\\
v) $\Leftrightarrow$ vi) 
Assume v) holds and let $\gamma$ be a positive map. By assumption, the composition $\gamma \circ \mathcal{E}$ is completely positive. Taking the adjoint preserves complete positivity, so that $
    (\gamma \circ \mathcal{E})^* = \mathcal{E}^* \circ \gamma^* $. Since $\gamma$ is positive, its adjoint $\gamma^*$ is also positive. Thus, for any positive map $\gamma^*$, the composition $\mathcal{E}^* \circ \gamma^*$ is completely positive. Because a channel is entanglement-breaking if and only if its adjoint is entanglement-breaking (as transposition preserves the separability of the Choi matrix), applying this reasoning to $\mathcal{E}$ directly yields vi).

\end{proof}
Point iii) of the theorem above implies that one only needs to check for the dimension $d_B$ in the identity $\id_d$ in the definition \ref{def:eb} of entanglement-breaking channels. An example of an entanglement-breaking channel is $\m E \in C_d, \rho \mapsto \trr{\rho}\1_{d'}/d'$ since its Choi matrix is $C_{\m E} = \1_d \otimes \1_{d'}$.
\begin{ex}\label{ex:depol-eb}
    The depolarizing channel $\Phi_p$ is entanglement-breaking if and only if $p\geq d/(d+1)$.
\end{ex}
\begin{proof}
The Choi matrix of the depolarizing channel is isotropic, for which separability is equivalent to the PPT condition. Parameterizing the channel as $\Phi_p(\rho) = (1-p)\rho + p \text{Tr}[\rho]\frac{\1}{d}$, its unnormalized Choi matrix evaluates to:
\begin{align*}
    C_{\Phi_p} &= (1-p)\ket{\Omega}\bra{\Omega} + p \left( \1 \otimes \frac{\1}{d} \right)
\end{align*}
Applying the partial transpose transforms $\ket{\Omega}\bra{\Omega}$ into the swap operator $\mathbb{F}$. The minimum eigenvalue of the swap operator is strictly $-1$ (occurring on the antisymmetric subspace). Enforcing the PPT condition dictates that the minimum eigenvalue of the fully transposed Choi matrix must remain non-negative:
\begin{align*}
    -(1-p) + \frac{p}{d} \ge 0 \Leftrightarrow
    p \left(1 + \frac{1}{d}\right) \ge 1 \Leftrightarrow
    p \ge \frac{d}{d+1}
\end{align*}
This directly yields the exact parameter regime where the channel becomes entanglement-breaking.
\end{proof}

The Choi matrix of a channel may have a positive partial transpose, but not necessarily be separable. The corresponding channels are called completely copositive or also PPT channels sometimes:
\begin{definition}[Complete copositivity/PPT channel]
    Let $\m E \in \m C(\C^d,\C^{d'}) $ be a quantum channel. If its Choi matrix is PPT, $(\id_d \otimes \vartheta_{d'})(C_T)\geq 0$, then the channel is called completely copositive. The set of all completely copositive channels will be denoted by PPT$(\C^d,\C^{d'})$ and PPT$_d\coloneqq \text{PPT}(\C^d)\coloneqq \text{PPT}(\C^d,\C^{d})$.
\end{definition}
This is equivalent to $(\id \otimes \m E)(\rho)$ being a PPT state again, for every input.
\begin{prop}
    The set of entanglement-breaking channels EB$_d$ is closed under composition with other entanglement-breaking quantum channels, as well as with other arbitrary channels, making it an ideal in the semigroup of quantum channels. The same is true for the set of PPT channels.
\end{prop}
\begin{proof}Let $\rho \in \m D(\C^d)$ be an aribrary state, then
    $(\id_d \otimes \Phi \circ \Psi)(\rho) =(\id_d \otimes \Phi)((\id_d \otimes \Psi)(\rho))$ is separable since $\Phi$ is entanglement-breaking. On the other hand, $(\id_d \otimes \Psi \circ \Phi)(\rho)$ is separable since product channels like $\id_d \otimes \Psi$ preserve separability. Similar reasoning proves the statement about PPT channels.
\end{proof}

The following conjecture was stated in \cite{christandl2012ppt}:
\begin{conjecture}[PPT$^2$ conjecture]. For any PPT quantum channel $\m E \in \text{PPT}_d$, its square is entanglement-breaking: $\m E^2 \in \text{EB}_d$.
\end{conjecture}
The PPT$^n$ conjecture is accordingly about $n$-time compositions. Of course, the trace-preservation of $\m E$ is irrelevant for the conjecture, so that it may be equivalently formulated for completely positive maps.



There is the following simple necessary channel for a channel to be entanglement-breaking:
\begin{lemma}\label{lem:rank}
    Let $\m E \in \m C_d$ be a quantum channel. If $\rk(C_{\m E}) < d$, then $C_T$ is entangled.
\end{lemma}
\begin{proof}
    Note that $C_{\m E}^\text{A} = \tr_\text{B}[C_{\m E}] = \1$, so that $\rk(C_{\m E}^\text{A}) = d$. Thus by proposition \ref{prop:rank-condition} it is $\rk(C_T) > d$.
\end{proof}
\begin{prop}[Necessary Kraus rank, theorem 6 in \cite{horodecki1999general}]
    Let $\m E \in \m C_d$. If $\m E$ can be written with fewer than $d$ Kraus operators, i.e., $\rk(C_{\m E})<d$, then $\m E$ is not entanglement-breaking.
\end{prop}
\begin{proof}
    The minimal number of required Kraus operators is the rank of the Choi matrix. Now use lemma \ref{lem:rank}.
\end{proof}
Unfortunately, there is no simple closed formula for the rank of the Choi matrix of a channel composition. Thus, this result does not provide an easy condition for the PPT$^n$ conjecture, for example.
\begin{prop}
    Let $\m E \in \m C_d$ be a quantum channel. If $\rk(C_{\m E}) = d$, then the channel $\m E$ is entanglement-breaking if and only if it is PPT.
\end{prop}
\begin{proof}
 Using proposition \ref{prop:low-rank} for the Choi matrix of the channel yields the statement.
\end{proof}

\subsubsection{Entanglement-Breaking Indices and Times}
Since entanglement-breaking channels are closed under composition with other quantum channels, in particular themselves, it may be interesting to investigate the smallest number of compositions needed for a channel to become entanglement-breaking (if that is the case). This concept was first investigated in \cite{lamiEntanglementBreakingIndices2015}. To take powers of a quantum channel, the channel must clearly be mapping between the same Hilbert space.
\begin{definition}[Entanglement-breaking index and time \cite{lamiEntanglementBreakingIndices2015, lamiEntanglementSavingChannels2016}]
    Let $\m E \in \m C_d$ be a quantum channel and $\{\Phi_t \in \m C_d\}_{t\geq 0}$ a quantum dynamical semigroup. The entanglement-breaking index of the channel and the entanglement-breaking time of the QDS are 
    \begin{align*}
        n_\text{EB}(T) &\coloneqq \min\{n_0 \in \N: \forall n \in \N, n\geq n_0 :\Phi^n \in \text{EB}_d\} \text{ and }\\
        t_\text{EB}(\Phi)&\coloneqq \inf\{t\geq 0: \forall s \geq t : \Phi_s \in \text{EB}_d \},
    \end{align*}
    respectively. If the above sets are empty, the respective entanglement-breaking index or time is set to $\infty$. Channels (QDSs) with a finite entanglement-breaking index (time) are called eventually entanglement-breaking. The set of eventually entanglement-breaking channels $\m E\in \m C_d$ will be denoted by EEB$_d$. If a channel is not eventually entanglement-breaking, it is entanglement-saving: $\text{ES}_d\coloneqq \m C_d $\textbackslash $\text{EEB}_d$
    \end{definition}\noindent
Since $t\mapsto \Phi_t$ is (strongly) continuous and EB$_d$ is closed, the entanglement-breaking time is achieved as a minimum. The entanglement-breaking indices of the elements of a QDS are closely linked to its entanglement-breaking time:
\begin{prop}
Let $(\Phi_t)_t$ be a QDS with entanglement-breaking time $t_\Phi$. Let $\delta \geq 0$, then the entanglement-breaking index of 
$\Phi_\delta$ is $n(\Phi_\delta) \in [\lfloor t_\Phi / \delta \rfloor, \lceil t_\Phi / \delta \rceil]$.
\end{prop}
\begin{proof}

Let $\Phi_\delta^n$ be entanglement-breaking, then $n\delta \geq t_\text{EB}(\Phi)$ and thus $ n \geq \lceil {t_\Phi}/{\delta} \rceil$ since $n \in \mathbb{N}$, which provides the desired upper bound for the entanglement-breaking index $n_\text{EB}(\Phi_\delta)$. On the other hand, clearly, 
$t_\text{EB}(\Phi) \geq \lfloor \frac{t_\text{EB}(\Phi)}{\delta} \rfloor \cdot \delta$ holds true.  Consequently, $\Phi_\delta^{\lfloor t_\Phi / \delta \rfloor}$ is entanglement-breaking if and only if ${t_\Phi}/{\delta} \in \mathbb{N} \Rightarrow n_0 \geq \lfloor \frac{t_\Phi}{\delta} \rfloor$.
\end{proof}\noindent
A QDS is, in particular, eventually entanglement-breaking if and only if any and thus all of its elements are.

\begin{ex} The entanglement-breaking index of the depolarizing map $\Phi_p$ is $n(\Phi_p) = \left\lceil - {\ln(d+1)}/{\ln(1-p)} \right\rceil$ for all $p \in\ ]0, 1[\ $, as well as $n(\Phi_0) = \infty$ and $n(\Phi_1) = 1$. 
\end{ex}
\begin{proof}Use example \ref{ex:depol-eb} and assume $p \in\ ]0, 1[\ $:
\begin{align*}
1 - (1-p)^n \ge \frac{d}{d+1} &\Leftrightarrow (1-p)^n \le 1 - \frac{d}{d+1} = \frac{d+1-d}{d+1} = \frac{1}{d+1} \\
&\Leftrightarrow n \ge \log_{1-p} \left(\frac{1}{d+1}\right) = \frac{\ln\left(\frac{1}{d+1}\right)}{\ln(1-p)} = -\frac{\ln(d+1)}{\ln(1-p)}.
\end{align*}
\end{proof}


If a channel $\m E \in \m C_d$ is not eventually entanglement-breaking, one might still consider the accumulation points of the sequence $(\m E^n)_{n\in \N}$.
\begin{definition}[AES channels \cite{lamiEntanglementSavingChannels2016} definition 3]
    Let $\m E \in \m C_d$ be a quantum channel. If all accumulation\footnote{Note that in \cite{lamiEntanglementSavingChannels2016} accumulation points are called limit points.} points of the sequence $(\m E^n)_{n \in \N}$ are entanglement-breaking, then $\m E$ is asymptotically entanglement-breaking. Otherwise, the channel is called asymptotically entanglement-saving. The set of asymptotically entanglement-breaking and entanglement-saving channels will be denoted by AEB$_d$ and AES$_d$, respectively.
\end{definition}
Interestingly, it is proven in proposition 29 of \cite{lamiEntanglementSavingChannels2016} that either all or none of the accumulation points are entanglement-breaking. The broad proof strategy is as follows: i) the set of all accumulation points forms a compact group, i.e., they are in particular invertible ii) by inverting the entanglement-breaking accumulation point, all other points can be written as a composition of this particular entanglement-breaking accumulation point and another auxiliary channel iii) use the ideal property of entanglement-breaking channels.

Let it be noted that in \cite{lamiEntanglementSavingChannels2016} they also introduce the notion of universal entanglement-preserving channels (UEP), which are those channels $\m E$ for which $(\id \otimes \m E)(\rho)$ is entangled whenever $\rho$ is. Such UEP channels are shown to be equivalent to the unitary channels. Unitary channels are addiionally entanglement-saving.

\subsubsection{Primitive channels}
Primitive quantum channels provide an example of channels that are eventually entanglement-breaking. This is a well-known result, see e.g. theorem 2.1 \cite{hansonEventuallyEntanglementBreaking2020} for a stronger result, where one direction is proven similarly to the present work.

If $\m E \in C_d$ is a primitive quantum channel, then by definition \ref{def:primitive-etc} it is 
\begin{align}
    \tilde{ \m E }\coloneqq \lim_{n \rightarrow \infty}\m E^n = \m R_\sigma \text{ where } \m D(\C^d)\ni\sigma > 0.
\end{align}
Hence, $C_{\tilde {\m E}}= \sum_{ij}\ket{i} \bra j \otimes \delta_{ij}\sigma = \1 \otimes \sigma$ is a product operator, so the limit $\tilde  {\m E}$ is an entanglement-breaking channel. In principle, the proof that primitive channels are eventually entanglement-breaking works by showing that the limit lies in the interior of EB$_d$ so that there must exist finite powers that come arbitrarily close to $\tilde{\m E}$ and thus also lie in $EB_d$.
\begin{lemma}
    A product operator $O = O_A \otimes O_B$ is in the interior of the set of separable operators if and only if $O_A$ and $O_B$ have full-rank:
    \begin{align*}
        O_A \otimes O_B \in \inter(\sep_{d_A d_B}) \Leftrightarrow (O_A > 0 \wedge O_B > 0)
    \end{align*}
\end{lemma}
\begin{proof}
One needs to show that there exists $\varepsilon > 0 \text{ such that } B_\varepsilon(O_A \otimes O_B) \subset \text{SEP}$. Obviously $\text{int}(\text{SEP}) \subset \text{int}(\mathcal{D}(\mathcal{H})) = \{ \text{full rank states} \}$. Hence, if $O_A \otimes O_B \in \text{int}(\text{SEP})$, then in particular $
    O_A \otimes O_B > 0$. Given the spectral decompositions $O_A = \sum_i \lambda_i |\lambda_i\rangle\langle\lambda_i|$ and $O_B = \sum_j \mu_j |\mu_j\rangle\langle\mu_j|$, one has 
$O_A \otimes O_B = \sum_{i,j} \lambda_i \mu_j (|\lambda_i\rangle \otimes |\mu_j\rangle)(\langle\lambda_i| \otimes \langle\mu_j|) > 0$, so that $O_A,O_B > 0$ follows from $O_A \otimes O_B > 0$. Hence, $(\Rightarrow)$ is done. Now assume that $O_A, O_B > 0$. Clearly, the identity on a bipartite system is a product matrix. Define the map
\begin{align*}
    \Lambda : \mathcal{B}(\mathcal{H}_{AB}) \to \mathcal{B}(\mathcal{H}_{AB}), \quad X \mapsto (O_A \otimes O_B)^{1/2} X (O_A \otimes O_B)^{1/2},
\end{align*}
such that $\Lambda(\1) = O_A \otimes O_B = O$. Consider the restriction $\Lambda|_{\text{SEP}}$; then $\text{Im}(\Lambda|_{\text{SEP}}) \subset \text{SEP}$ because the map is linear and preserves tensor products. By theorem 1 in \cite{gurvitsLargestSeparableBalls2002}, the identity matrix $\1$ lies in the interior of the separable cone, meaning there exists an open neighborhood $U \subset \text{SEP}$ containing $\1$. Furthermore, because $O_A, O_B > 0$, their square roots are uniquely defined and strictly invertible. This makes $\Lambda$ an invertible linear map on a finite-dimensional vector space, which guarantees that it is an open map. Consequently, $\Lambda$ maps the open set $U$ to an open set $\Lambda(U)$ containing $\Lambda(\1) = O$. Since $\Lambda(U) \subset \text{SEP}$, there exists an open neighborhood around $O$ that is entirely contained within the separable set, proving that $O \in \text{int}(\text{SEP})$.

\begin{prop}
    Primitive channels are eventually entanglement-breaking. \label{prop:primitive-eeb}
\end{prop}
\begin{proof}
    Since the Choi matrix $C_{\tilde{\m E}}$ of the limit of a primitive channel $\m E$ is a product matrix, it lies in $\inter(\sep_d)$, which is an open subset of $\sep_d$. Thus, there is an $\epsilon_{\m E} > 0$, such that $B_{\epsilon} (C_{\tilde {\m E}})\in \sep_d$. By the definition of convergence, there is $N(\epsilon) \in \N$ such that for all $n> N$ it is $\norm{\m E^n- \tilde{\m E}}< \epsilon$. The Choi mapping $\m C: \m E \mapsto C_{\m E}$ is linear map, thus bounded continuous. Consequently, one has  $\norm{C_{\m E}-C_{\tilde{\m E}}} = \norm{\m C({\m E}-{\tilde{\m E}})} \leq \norm {\m C} \norm{{\m E}-{\tilde{\m E}}}$. Now, for all $n> N(\epsilon_{\m E}/\norm {\m C})$ it holds that $\norm{C_{\m E^n}-C_{\tilde{\m E}}} < \epsilon_{\m E}$ and thus $C_{\m E^n} \in \sep_d$ and $\m E^n \in EB_d$.
\end{proof}

One potentially useful observation is that to prove the $\text{PPT}^2$ conjecture, one may restrict to a dense set:

\begin{prop}[Corollary 3.2 in \cite{hansonEventuallyEntanglementBreaking2020}] The $\text{PPT}^n$ conjecture holds if and only if every PPT primitive quantum channel $\m E$ has $\m E^n \in \text{EB}_d$.
\end{prop}
\begin{proof} Let $\m E \in \text{PPT}_d$ be a PPT quantum channel. By denseness of primitive quantum channels, there is a sequence $\m E_k \rightarrow \m E$ with each $\m E_k$ being PPT and primitive. By continuity of compositions, it is $\m E^n = \lim_{k \to \infty} \m E_k^n$. Since the set of entanglement-breaking channels is closed, $\m E^n \in \text{EB}_d$.
\end{proof}
By reading the proof, one readily sees that the statement holds for any dense set.

\subsubsection{Other results in the literature}
In the case of quantum dynamical semigroups, also the converse of proposition \ref{prop:primitive-eeb} holds, that is, every eventually entanglement-breaking QDS is primitive. This is shown in theorem 2.1 of \cite{lamiEntanglementSavingChannels2016}. This paper also goes further on to provide tighter bounds for the entanglement-breaking times of quantum dynamical semigroups.

There are many more recent developments on the PPT conjecture.  Notably, Rahaman et al. \cite{rahamanEventuallyEntanglementBreaking2018} rigorously demonstrated that every PPT channel is eventually entanglement-breaking (Theorem 4.4). Their work further characterizes the asymptotic behavior of unital maps: While the identity map serves as a trivial counterexample preventing a universal claim, they proved that almost all unital channels are indeed entanglement-breaking (Theorem 5.2). Additionally, they established that any unital irreducible channel is asymptotically entanglement-breaking (Corollary 6.1).  In \cite{christandlWhenComposedMaps2019} proved the conjecture for dimension $d=3$ using Schmidt number techniques and formulated several equivalent conjectures that provide alternative avenues for understanding this behavior.



  
